\subsection{Selected vectors with different sector probabilities}
\label{sec:ldp_sector_weights}

The logarithmic-depth conversion remark in the discussion and outlook
of~\cite{Malz2024Preparation} leaves the choice of vector in a degenerate
parent groundspace unspecified. Equality of those groundspaces alone does
not imply a shallow unitary conversion between the selected vectors.
The following obstruction concerns this interpretation; it does not impose
a definition of phase or exclude conversion of compatible coherent seeds.

\begin{theorem}[Matrix elements outside the support]
    \label{thm:ldp_offsupport_matrix_elements}
    \lean{QuantumCircuit.apply_eq_zero_of_mem_supportedOperators,
        QuantumCircuit.productVector_single,
        QuantumCircuit.expect_single,
        QuantumCircuit.mul_apply_eq_mul_of_mem_supportedOperators,
        QuantumCircuit.expect_smul_single_add_of_offDiag_zero}
    \leanok
    \uses{def:qc_local_circuit}
    Let $A$ act on $S$. If configurations $\sigma,\tau$ differ at a site
    outside $S$, then $\langle\sigma|A|\tau\rangle=0$.
    If $A,B$ act on disjoint sets, then
    $\langle\sigma|AB|\sigma\rangle=
        \langle\sigma|A|\sigma\rangle\langle\sigma|B|\sigma\rangle$.
\end{theorem}
\begin{proof}
    \leanok
    A product operator with identity factors outside $S$ has an entry
    containing $\delta_{\sigma_k,\tau_k}=0$ at the differing site $k$.
    Extend this equality by linearity. A computational-basis vector is
    the product of its one-site basis vectors, so disjoint expectations
    factor.
\end{proof}

\begin{theorem}[Stability of connected correlations]
    \label{thm:ldp_covariance_stability}
    \lean{QuantumCircuit.expect_eq_inner,
        QuantumCircuit.norm_expect_le_one,
        QuantumCircuit.norm_expect_sub_le,
        QuantumCircuit.expect_smul_state,
        QuantumCircuit.norm_covariance_sub_le}
    \leanok
    For unit vectors $\psi,\phi$ and unitary observables $A,B$, write
    $\operatorname{Cov}_\psi(A,B)=
        \langle\psi|AB|\psi\rangle-
        \langle\psi|A|\psi\rangle\langle\psi|B|\psi\rangle$.
    Then $|\langle\psi|A|\psi\rangle|\le1$, and
    \begin{align}
        |\langle\psi|A|\psi\rangle-
        \langle\phi|A|\phi\rangle| & \le2\|\psi-\phi\|, \\
        |\operatorname{Cov}_\psi(A,B)-\operatorname{Cov}_\phi(A,B)|
                                   & \le6\|\psi-\phi\|.
    \end{align}
    Both expectations and covariances are unchanged by a unit scalar
    multiplying the state.
\end{theorem}
\begin{proof}
    \leanok
    Expand the expectation difference as
    $\langle\psi-\phi|A\psi\rangle+
        \langle\phi|A(\psi-\phi)\rangle$ and apply Cauchy--Schwarz.
    Expand the difference of the two products of expectations into two
    terms. The three expectation differences each contribute at most
    $2\|\psi-\phi\|$.
\end{proof}

\begin{theorem}[Separated correlations in a weighted GHZ vector]
    \label{thm:ldp_weighted_ghz_covariance}
    \lean{QuantumCircuit.ghzState_two_eq,
        QuantumCircuit.zero_config_ne_one,
        QuantumCircuit.norm_ghzState_two,
        QuantumCircuit.norm_ghzState_two_sq,
        QuantumCircuit.ghzState_smul_two,
        QuantumCircuit.expect_ghzState_two_diagonal,
        QuantumCircuit.expect_ghzState_two,
        QuantumCircuit.covariance_ghzState_two,
        QuantumCircuit.norm_covariance_ghzState_two_le,
        QuantumCircuit.norm_covariance_mulVec_ghzState_two_le,
        QuantumCircuit.diagonal_comp_eval_mem_supportedOperators}
    \leanok
    Let $|a|^2+|b|^2=1$, $N\ge1$, and
    $\chi=a|0^N\rangle+b|1^N\rangle$.
    Let unitary observables $X,Y$ act on disjoint sets $I,J$ with
    nonempty complement. Put $x_s=\langle s^N|X|s^N\rangle$ and
    $y_s=\langle s^N|Y|s^N\rangle$. For arbitrary $a,b$, $\|\chi\|^2=|a|^2+|b|^2$, and scaling $\chi$
    by $c$ scales both coefficients by $c$. Under the normalization above,
    $\|\chi\|=1$ and
    \begin{align}
        \operatorname{Cov}_\chi(X,Y)
                                       & =|a|^2|b|^2(x_0-x_1)(y_0-y_1), \\
        |\operatorname{Cov}_\chi(X,Y)| & \le4|a|^2|b|^2.
    \end{align}
    The same bound holds after a depth-$T$ circuit for observables whose
    backward light cones are disjoint and leave a site unobserved.
\end{theorem}
\begin{proof}
    \leanok
    \uses{thm:ldp_offsupport_matrix_elements,thm:ldp_covariance_stability}
    The unobserved site kills both cross terms. Within each product branch,
    $\langle XY\rangle_s=x_sy_s$. Subtract the product of the two weighted
    expectations and use $|a|^2+|b|^2=1$. Each diagonal expectation has
    modulus at most one. For a circuit, pull both observables back through
    the unitary, retaining their support inside the backward light cones.
\end{proof}

% TENKZ-SECTOR-WEIGHT-BEGIN
% Formula: <s^N| X_I Y_J |t^N> = <s_I|X|t_I><s_J|Y|t_J><s_R|t_R>,
% with R the nonempty complement of I union J. For s != t the last factor is zero.
% Ink: columns are the I, J, R Hilbert spaces; all six wires are physical contractions.
% No open legs or bonds between columns: the boundary signature is the empty scalar signature.
\begin{center}
    \tenkzkernel
    \begin{tenkz}[rows={wire,wire,wire}, cols=3, bonds=none]
        \tn[at=(1,1), name=brai, skin=box, ports={270:physical}]{\langle s_I\rvert}
        \tn[at=(1,2), name=braj, skin=box, ports={270:physical}]{\langle s_J\rvert}
        \tn[at=(1,3), name=brar, skin=box, ports={270:physical}]{\langle s_R\rvert}
        \tn[at=(2,1), name=obsx, skin=box, ports={90:physical,270:physical}]{X}
        \tn[at=(2,2), name=obsy, skin=box, ports={90:physical,270:physical}]{Y}
        \tn[at=(2,3), name=obsr, skin=box, ports={90:physical,270:physical}]{1_R}
        \tn[at=(3,1), name=keti, skin=box, ports={90:physical}]{\lvert t_I\rangle}
        \tn[at=(3,2), name=ketj, skin=box, ports={90:physical}]{\lvert t_J\rangle}
        \tn[at=(3,3), name=ketr, skin=box, ports={90:physical}]{\lvert t_R\rangle}
        \tnwire{brai.270}{obsx.90}
        \tnwire{braj.270}{obsy.90}
        \tnwire{brar.270}{obsr.90}
        \tnwire{obsx.270}{keti.90}
        \tnwire{obsy.270}{ketj.90}
        \tnwire{obsr.270}{ketr.90}
    \end{tenkz}
\end{center}
% TENKZ-SECTOR-WEIGHT-END

\begin{theorem}[Equal parent Hamiltonians with unequal periodic weights]
    \label{thm:ldp_sector_weight_parent}
    \lean{MPSTensor.groundSpaceMap_repeatedBlockTensor_eq_ghzState,
        MPSTensor.groundSpaceMap_ghzTensor_eq_ghzState,
        MPSTensor.mpv_repeatedBlockTensor_eq_ghzState,
        MPSTensor.mpv_ghzTensor_eq_ghzState,
        MPSTensor.norm_mpvState_repeatedBlockTensor,
        MPSTensor.norm_mpvState_ghzTensor,
        MPSTensor.normalizedMPVState_repeatedBlockTensor,
        MPSTensor.normalizedMPVState_ghzTensor,
        MPSTensor.groundSpace_repeatedBlockTensor_eq_ghzTensor,
        MPSTensor.parentInteraction_repeatedBlockTensor_eq_ghzTensor,
        MPSTensor.chainGroundSpace_repeatedBlockTensor_eq_ghzTensor,
        MPSTensor.parentHamiltonian_repeatedBlockTensor_eq_ghzTensor}
    \leanok
    Define
    \begin{align}
        A^0 & =\operatorname{diag}(1,1,0), & A^1 & =\operatorname{diag}(0,0,1), \\
        B^0 & =\operatorname{diag}(1,0),   & B^1 & =\operatorname{diag}(0,1).
    \end{align}
    On every positive chain, their normalized periodic vectors are
    $(2|0^N\rangle+|1^N\rangle)/\sqrt5$ and
    $(|0^N\rangle+|1^N\rangle)/\sqrt2$.
    At every interaction length the local boundary-generated spaces and
    canonical parent projectors agree. Consequently their periodic-chain
    constraint spaces and canonical parent Hamiltonians agree for every
    chain length.
\end{theorem}
\begin{proof}
    \leanok
    A boundary matrix $X$ for $A$ contributes amplitudes
    $(X_{00}+X_{11},X_{22})$; one for $B$ contributes $(X_{00},X_{11})$.
    Both pairs range over $\mathbb C^2$. For the identity boundary they
    become $(2,1)$ and $(1,1)$, whose norms on a positive chain are
    $\sqrt5$ and $\sqrt2$. The local orthogonal projectors are determined
    by these boundary-generated spaces. The same local terms define the
    chain constraints and Hamiltonians.
\end{proof}

\begin{theorem}[A lower bound for changing sector probabilities]
    \label{thm:ldp_sector_weight_obstruction}
    \lean{MPSPreparation.ghz_sector_weight_infidelity_lower_bound,
        MPSTensor.repeatedBlockTensor_infidelity_lower_bound,
        MPSTensor.length_le_of_repeatedBlockTensor_infidelity_lt,
        QuantumCircuit.isSeparatedBy_of_val_bounds,
        QuantumCircuit.IsSeparatedBy.symm,
        QuantumCircuit.IsSeparatedBy.notMem_neighbourhood,
        QuantumCircuit.exists_separated_singletons_notMem_neighbourhood}
    \leanok
    \uses{def:qc_local_circuit}
    Let $\psi_N=(2|0^N\rangle+|1^N\rangle)/\sqrt5$ and
    $\phi_N=(|0^N\rangle+|1^N\rangle)/\sqrt2$.
    Every depth-$T$ nearest-neighbour unitary circuit $U$ on the ring,
    with $N>4T+4$, satisfies
    \begin{align}
        1-|\langle U\psi_N|\phi_N\rangle|\ge\frac9{5000}.
    \end{align}
    These are the normalized periodic vectors of the tensors in
    Theorem~\ref{thm:ldp_sector_weight_parent}. Conversely, overlap error
    strictly below $9/5000$ forces $N\le4T+4$.
\end{theorem}
\begin{proof}
    \leanok
    \uses{thm:ldp_weighted_ghz_covariance,thm:ldp_covariance_stability,
        lem:ldp_error_triangle}
    Choose sites $i=0$, $j=2T+2$ and $k=T+1$. The radius-$T$ backward
    cones of $i,j$ are disjoint, and neither contains $k$, including around
    the periodic boundary. The Pauli observables $Z_i,Z_j$ have covariance
    at most $16/25$ in modulus in $U\psi_N$, while their covariance in
    $\phi_N$ is one. Choose $|c|=1$ with
    $\eta^2=\|U\psi_N-c\phi_N\|^2
        =2(1-|\langle U\psi_N|\phi_N\rangle|)$.
    Correlation stability gives $1\le16/25+6\eta$, hence
    $\eta\ge3/50$ and the stated bound.
\end{proof}
